% =========================================================================
%  CHAPTER 4: THE THREE SCREENING REGIMES
%  Volume I: The Evidence — Part B
% =========================================================================
%  STATUS: COMPLETE
%  SOURCE: Extracted from Book_3/Chapter_11 §6 (screening analysis)
% =========================================================================

\chapter{The Three Screening Regimes}
\label{ch:screening-regimes}

The previous chapter proved that $\kop = 0.5464$ is universal:\ it does not change with electron count, nuclear charge, or the complexity of the atomic system.  All multi-electron complexity lives in a single function: the \textbf{screening constant} $\sigma(Z, N)$.

This chapter maps the structure of that function.


% =============================================
\section{The Per-Electron Screening Efficiency}
\label{sec:per-electron}
% =============================================

The raw screening constants span a vast range --- from $\sigma = 0$ (hydrogen) to $\sigma = 74$ (gold).  To extract the physically meaningful signal, we normalise by the number of screening electrons:

\begin{equation}\label{eq:efficiency}
  \eta \equiv \frac{\sigma}{N - 1}
\end{equation}

This is the \textbf{per-electron screening efficiency}: the average fraction of the nuclear charge occluded by each inner electron.  Its evolution across the periodic table reveals the geometric structure of electron shells:

\begin{center}
\begin{tabular}{rlrcl}
\toprule
$N$ & \textbf{Sequence} & $\eta = \sigma/(N{-}1)$ & \textbf{Shell type} & \textbf{Regime} \\
\midrule
  2 & He-like  & 0.620 & 1s$^2$           & \textbf{I} \\
  3 & Li-like  & 0.812 & [He]\,2s          & \textbf{II} \\
 10 & Ne-like  & 0.781 & [He]\,2s$^2$2p$^6$ & \textbf{II} \\
 18 & Ar-like  & 0.821 & [Ne]\,3s$^2$3p$^6$ & \textbf{II} \\
 28 & Ni-like  & 0.922 & [Ar]\,3d$^{10}$4s$^0$ & \textbf{III} \\
 46 & Pd-like  & 0.935 & [Kr]\,4d$^{10}$    & \textbf{III} \\
 79 & Au-like  & 0.931 & [Xe]\,4f$^{14}$5d$^{10}$6s$^1$ & \textbf{III} \\
\bottomrule
\end{tabular}
\end{center}

Three distinct regimes emerge.  They are not arbitrary --- each corresponds to a qualitatively different geometric arrangement of inner electrons relative to the nucleus and the outermost electron.


% =============================================
\section{Regime I: Dyad Occlusion ($\eta \approx 0.62$)}
\label{sec:regime1}
% =============================================

\textbf{System:} Helium-like ions ($N = 2$, both electrons in the 1s shell).

In this regime, two electrons share the \emph{same} shell.  Neither is ``inside'' the other.  The screening is not radial (one shell shielding the next) but \textbf{angular}: each electron occludes a cone of the nuclear pressure field as seen from the position of the other.

\subsection{The Geometry}

Two 1s electrons are diametrically opposed on the nuclear surface (anti-parallel spins, maximum separation).  Each subtends a solid angle as seen from the other:
\begin{equation}
  \Omega_{\text{e}} = 2\pi\left(1 - \cos\theta_{\text{e}}\right)
\end{equation}

The screening efficiency $\eta = 0.62$ implies that each electron occludes $62\%$ of the nuclear field from its partner.  This is a \emph{geometric floor}: it is impossible for a same-shell electron to occlude more than $\sim 65\%$ of the nuclear field, because it subtends less than $2\pi$ steradians from the partner's viewpoint.

\subsection{Z-Dependence}

The screening in this regime shows a subtle $Z$-dependence, declining from $\sigma = 0.656$ (He) to $\sigma = 0.527$ (Fe$^{24+}$).  At high $Z$, both electrons are compressed closer to the nucleus, their angular separation decreases, and the occlusion cone narrows.

\textbf{SDT interpretation:} Pressure shadow compression.  As the nuclear pressure increases, the electron vortices are squeezed into a tighter solid angle, reducing their mutual geometric occlusion.


% =============================================
\section{Regime II: Shell-Layered Occlusion ($\eta \approx 0.78\text{--}0.82$)}
\label{sec:regime2}
% =============================================

\textbf{Systems:} Lithium-like through argon-like ($N = 3$ to $18$, shells $n = 2$ to $3$).

In this regime, the outermost electron occupies a \emph{different shell} from the core electrons.  The inner electrons are \emph{between} the nucleus and the valence electron, creating radial shielding.

\subsection{Why Screening Is More Efficient}

Radial shielding is geometrically more efficient than angular shielding.  An inner-shell electron directly intercepts the nuclear pressure field that would otherwise reach the outer electron.  The $62\%$ of Regime~I rises to $78\text{--}82\%$.

\subsection{The Saturation Effect}

Within Regime~II, $\eta$ shows a characteristic pattern:
\begin{itemize}
  \item $N = 3$: $\eta = 0.812$ (2 core electrons shielding 1 valence)
  \item $N = 10$: $\eta = 0.781$ (9 electrons shielding 1 valence)
  \item $N = 18$: $\eta = 0.821$ (17 electrons shielding 1 valence)
\end{itemize}

The slight \emph{dip} at $N = 10$ is significant.  When the $n = 2$ shell is full (2s$^2$\,2p$^6$), the valence electron (2p$^6$) is at the \emph{outer edge} of its own shell.  Its peer electrons within the same shell provide angular screening (Regime~I-like), pulling $\eta$ down slightly.

At $N = 18$, the valence electron is in $n = 3$, and the \emph{entire} $n = 2$ shell is now radially interior, restoring high radial screening efficiency.

\subsection{Physical Mechanism}

In SDT language: s and p electrons create \textbf{layered pressure shadows}.  Each shell is a concentric barrier that intercepts a fraction of the nuclear field.  The cumulative effect is multiplicative but with diminishing returns: each layer shields part of what the previous layer already shielded.

The overlap of pressure shadows limits per-electron efficiency to $\sim 82\%$.  Breaking through this ceiling requires a different orbital geometry.


% =============================================
\section{Regime III: Geometric Lock Occlusion ($\eta \approx 0.92\text{--}0.95$)}
\label{sec:regime3}
% =============================================

\textbf{Systems:} Nickel-like through gold-like ($N = 28$ to $79$, including d and f shells).

The introduction of d-electrons produces a \textbf{discontinuous jump} in per-electron screening efficiency: from $\sim 0.82$ (Regime~II) to $\sim 0.92$ (Regime~III).  A $12\%$ increase that cannot be explained by gradual shell-layering.

\subsection{Why d-Electrons Are Different}

The d-orbitals have a fundamentally different geometry from s and p orbitals:
\begin{itemize}
  \item \textbf{s-orbitals:} spherically symmetric.  One lobe.
  \item \textbf{p-orbitals:} two lobes, oriented along one axis.  Three per shell.
  \item \textbf{d-orbitals:} four lobes, oriented between axes.  Five per shell.
\end{itemize}

Five d-orbitals together tile solid angle with far greater coverage than three p-orbitals.  A complete d$^{10}$ subshell (10 electrons in 5 orientations, each with 2 spin states) creates a nearly hermetic geometric enclosure around the nucleus.

\subsection{The Two-Halves Principle}

A half-filled d-shell (d$^5$) fills one geometric hemisphere.  A complete d-shell (d$^{10}$) fills both hemispheres, creating a \textbf{double-pentagon geometric lock}: a configuration where the five d-orbital orientations tile the sphere with minimal gaps.

This geometric lock explains:
\begin{enumerate}
  \item The anomalous stability of d$^5$ (Cr, Mn) and d$^{10}$ (Cu, Zn) configurations.
  \item The preference of palladium for [Kr]\,4d$^{10}$ with \emph{zero} 5s electrons.
  \item The exceptional catalytic properties of transition metals (partially unlocked d-shells provide geometric access to the nuclear field).
\end{enumerate}

\subsection{The f-Electron Ceiling}

Adding f-electrons ($N = 79$, gold-like) does \emph{not} significantly increase $\eta$ beyond the d-shell value:
\begin{itemize}
  \item $N = 28$ (Ni-like, first d$^{10}$): $\eta = 0.922$
  \item $N = 46$ (Pd-like, second d$^{10}$): $\eta = 0.935$
  \item $N = 79$ (Au-like, $+$ f$^{14}$): $\eta = 0.931$
\end{itemize}

The f-electrons are so deeply buried that their \emph{incremental} contribution to outer-shell screening is negligible.  They are already ``priced in'' to the inner shell compression.  The per-electron efficiency \textbf{saturates} near $0.93$, a geometric ceiling that no additional electrons can breach.


% =============================================
\section{The Discontinuity: Where Chemistry Changes}
\label{sec:discontinuity}
% =============================================

The jump from Regime~II to Regime~III at $N \approx 28$ is not a gradual transition.  It is a \textbf{phase boundary} in screening efficiency:

\begin{center}
\begin{tabular}{rl}
\toprule
$N$ & $\eta$ \\
\midrule
18 (Ar-like, last s/p-only) & 0.821 \\
28 (Ni-like, first d$^{10}$) & 0.922 \\
\midrule
\textbf{Jump:} & \textbf{+0.101 ($+12\%$)} \\
\bottomrule
\end{tabular}
\end{center}

This jump coincides precisely with:
\begin{itemize}
  \item The onset of transition metals in the periodic table.
  \item The dramatic change in chemical behaviour (variable oxidation states, coloured ions, catalytic activity).
  \item The appearance of metallic bonding and band structure.
\end{itemize}

The screening discontinuity is not merely a curiosity of atomic structure.  It is the \textbf{geometric origin of the periodic table's division} between main-group and transition-metal chemistry.


% =============================================
\section{Summary: The Screening Landscape}
\label{sec:screening-summary}
% =============================================

\begin{enumerate}
  \item All multi-electron complexity in atomic structure resides in $\sigma(Z, N)$, not in $\kop$.
  \item The per-electron efficiency $\eta = \sigma/(N{-}1)$ reveals three geometric regimes:
    \begin{itemize}
      \item \textbf{Regime I} ($\eta \approx 0.62$): Same-shell angular occlusion (the geometric floor).
      \item \textbf{Regime II} ($\eta \approx 0.78\text{--}0.82$): Cross-shell radial occlusion (layered pressure shadows).
      \item \textbf{Regime III} ($\eta \approx 0.92\text{--}0.95$): d/f-shell geometric lock (the geometric ceiling).
    \end{itemize}
  \item The $12\%$ jump at the d-shell boundary is the microscopic origin of the periodic table's division between s/p and d-block chemistry.
  \item The ceiling at $\eta \approx 0.93$ represents maximum geometric solid-angle coverage by electron vortices.
  \item These three regimes are the microscopic origin of the three coefficients ($k_{sp}$, $k_d$, $k_f$) in the Recursive Shell Compression Rule (next chapter).
\end{enumerate}

\subsection{Epistemic Status}

The screening efficiency values $\eta \approx 0.62$, $0.80$, and $0.93$ are \textbf{empirically observed} from ionisation energy data (NIST Atomic Spectra Database).  They are not yet derived from first principles.

The geometric narratives presented above---dyad occlusion, shell-layered pressure shadows, geometric lock---are \textbf{interpretive frameworks} consistent with SDT ontology.  They are plausible, physically motivated, and qualitatively correct.  But they are not derivations.

A genuine geometric derivation would compute the solid-angle occlusion fraction of inner electrons on the nuclear pressure field for each shell configuration, and predict:
\begin{equation}
  \sigma(Z, N) = \mathcal{F}_{\text{occlusion}}(Z, N, \text{shell geometry})
\end{equation}
from first principles, without reference to $E_I$ data.  That derivation is an \textbf{open problem} (Chapter~\ref{ch:open-problems}).

Until it is completed, the three regimes are best understood as \emph{robust empirical phenomenology with a compelling geometric interpretation}, not as closed derivations.

